% =====================================================================
%  Section IV. Strong-Coupling Approximation
% =====================================================================
\section{STRONG-COUPLING APPROXIMATION}\label{sec:strong}

The gauge field average $Z(P)$ which determines whether quarks are bound
was defined on a lattice in Sec.~\ref{sec:lattice}
(Eq.~\eqref{eq:3.13}).  There are two limits in which this average can
be calculated.  The most interesting limit is the strong-coupling limit
$g \to \infty$.  This is the limit which exhibits quark binding.  A
strong-coupling expansion will be derived in this section.  The
strong-coupling expansion will be the basis for a reformulation of the
gauge-field theory as a string model.  This will also be explained in
this section.

Consider specifically the numerator of Eq.~\eqref{eq:3.13}, to be
denoted $I_{\theta}(P)$:
\begin{equation}
I_{\theta}(P) = \prod_{n,\mu\nu}\int\ddth{\mu\nu}
  \exp\!\Bigl[i\sum_{P}(\pm)\thl{\mu\nu}
    +\frac{1}{2g^{2}}\sum_{n,\mu\nu}
      \bigl(e^{if_{\mu\nu}(n)}+e^{-if_{\mu\nu}(n)}\bigr)\Bigr].
\label{eq:4.1}
\end{equation}

Expanding the integrand in powers of $1/g^{2}$, the zeroth-order term is
\begin{equation}
I_{\theta}^{(0)}(P) = \prod_{n,\mu\nu}\int\ddth{\mu\nu}
  \exp\!\Bigl[i\sum_{P}(\pm)\thl{\mu\nu}\Bigr].
\label{eq:4.2}
\end{equation}
This term vanishes, since for any $\thl{\mu\nu}$ which appears in
$\sum_{P}$, there is an integral
$\int\ddth{\mu\nu}\exp(\pm i\thl{\mu\nu})$ which is zero.  Thus one must
seek higher-order terms in $g^{-2}$ which cancel the $\thl{\mu\nu}$ in
the line integral.  The first-order term is
\begin{equation}
I_{\theta}^{(1)}(P) = \prod_{n,\mu\nu}\int\ddth{\mu\nu}
  \exp\!\Bigl[i\sum_{P}(\pm)\thl{\mu\nu}+if_{\mu\nu}(l)\Bigr].
\label{eq:4.3}
\end{equation}
The quantity $f_{\mu\nu}(l)$ is itself a line integral of the gauge
field; it is the line integral around a square originating at the lattice
site $l$ of size $a$ (unit square).  The integral for $I_{\theta}^{(1)}(P)$
will vanish unless it is possible to find a unit square such that
$f_{\mu\nu}(l)$ cancels completely the line integral
$\sum_{P}(\pm)\thl{\mu\nu}$.  This is possible only if the path $P$ is
itself a unit square.  Otherwise the first-order term vanishes and one
must study the terms of order $g^{-2}$ or higher.

The term of order $g^{-2n}$ has the form
\begin{equation}
I_{\theta}^{(n)}(P) = \prod_{n,\mu\nu}\int\ddth{\mu\nu}
  \exp\!\Bigl[i\sum_{P}(\pm)\thl{\mu\nu}
    + i\sum_{m=1}^{n}f_{\mu\nu}(l_{m})\Bigr].
\label{eq:4.4}
\end{equation}
The only nonzero terms in this sum are those for which
\begin{equation}
\sum_{m}f_{\mu\nu}(l_{m}) = 0.
\label{eq:4.5}
\end{equation}
[See Eq.~\eqref{eq:3.8} for the definition of $f_{\mu\nu}$ in terms of
$\thl{\mu\nu}$.]  This equation can be understood geometrically.  Each
$f_{\mu\nu}$ corresponds to a square of size $a$ on the lattice.  For
this sum to vanish the set of squares defined by $f_{\mu\nu}$,
$f_{\mu\nu'}$, \ldots\ must combine to make a surface with boundary $P$.
(To be precise, each $f_{\mu\nu}$ corresponds to a line integral around a
square, and when these squares are joined to make a surface the line
integrals must cancel along all internal lines of the surface.  The line
integrals along the boundary $P$ of the surface must run in the opposite
direction to the original path $P$.)  See Fig.~\ref{fig:10}.

For a given path $P$ the lowest nonzero order in $I_{\theta}(P)$ is
determined by the minimal area $A$ enclosed by $P$, the area $A$ being
the area of any surface built of unit squares on the lattice with
boundary $P$.  Then $I_{\theta}(P) \sim (g^{-2})^{A/a^{2}}$ apart from a
numerical factor.

\begin{figure}[t]
  \centering
  \includegraphics[width=0.5\textwidth]{fig10}
  \refstepcounter{figure}%
  \caption*{Fig. 10. Filling of enclosed area of path $P$ by elementary
    squares.}
  \label{fig:10}
\end{figure}

This is the result promised in Sec.~\ref{sec:binding}:  The gauge-field
average $I_{\theta}(P)$ behaves as
$\exp[-A\,(\ln g^{2})/a^{2}]$, i.e., exponentially in the area enclosed
by $P$.  Hence, according to the arguments of Sec.~\ref{sec:binding},
quark paths will not separate macroscopically, and there will be no
quarks among the final-state particles.

Consider higher-order terms in the expansion of $I_{\theta}(P)$ for
given $P$.  There are many such terms because there are many surfaces
with boundary $P$.  In particular, there are many ways to combine subsets
of $f$'s to add to zero so such subsets can be added to any minimal sum
of $f$'s which forms a surface with boundary $P$.  The simplest example
of a set of $f$'s which add to zero are the set of $f$'s corresponding
to the six faces of a unit cube.  Written out, this gives
\begin{equation}
\begin{split}
&f_{\mu\nu}(n)+f_{\mu\nu}(n+\hat\mu)+f_{\mu\nu}(n+\hat\nu)\\
&\qquad + f_{\mu\rho}(n)+f_{\rho\mu}(n+\hat\mu)+f_{\mu\rho}(n+\hat\rho)=0,
\end{split}
\label{eq:4.6}
\end{equation}
which is easily checked using Eq.~\eqref{eq:3.8}.  [Equation
\eqref{eq:4.6} is the lattice analog of the equation
$\epsilon_{\mu\nu\rho\sigma}\partial_{\rho}F_{\mu\nu}(x)=0$.]

Let $A(P)$ be the minimal area as defined above enclosed by $P$.  Since
one can place a unit cube anywhere on this minimal area, it means that
there are roughly $A(P)/a^{2}$ more terms of order $g^{-2(n+1)}$ in the
expansion of $I_{\theta}(P)$ than there are terms of order $g^{-2n}$.
[One can place unit cubes anywhere in space, not just on the minimal
surface; but when one divides $I_{\theta}(P)$ by $Z$ all disconnected
terms cancel, as usual.]  This suggests that the $1/g^{2}$ expansion is
not very useful in the limit $A(P) \to \infty$, which is the limit of
interest for quark binding.  However, experience with related problems
suggests that $I_{\theta}(P)$ is not the appropriate quantity to expand;
instead one should try writing
\[
I(P) = Z^{-1}I_{\theta}(P) = (g^{-2})^{A(P)/a^{2}}e^{-c(P,g^{2})}
\]
and expand $c(P,g^{2})$ in powers of $g^{-2}$ instead.  One would expect
$c(P,g^{2})$ to be dominated by a term proportional to $A(P)$, say
\[
c(P,g^{2}) = A(P)f(g^{-2}) + \order{P}
\]
(where $P$ is the length of the path $P$).  The crucial question is the
nature of the series for $f(g^{-2})$.  Past experience with similar types
of expansions (namely, the high-temperature expansions of statistical
mechanics:  see, e.g., Ref.~13) suggests that $f(g^{-2})$ will have a
convergent expansion at least for $g^{-2}$ less than a critical value
$g_{c}^{-2}$.  However, no calculations have been done in the
gauge-field theory for $f(g^{-2})$ as yet.

Consider the complete expansion of $I(P)$.  Each nonzero term in the
expansion corresponds to a surface with perimeter $P$.  The complete
expansion corresponds to a sum over all possible surfaces with given
perimeter $P$.  ``All possible'' surfaces include surfaces which
intersect themselves (to take into account terms where a given
$f_{\mu\nu}$ appears several times in the sum
$\sum f_{\mu\nu} + f_{\mu\nu'} + \cdots$).  There is a weight factor for
each surface, aside from the power of $g^{-2}$ determined by the area of
the surface.  For a simple surface, the weight factor is 1; the weight is
more complicated for self-intersecting surfaces.

Thus, the strong-coupling expansion for the current-current propagator
has the same general structure as in string models of hadrons.  One is
actually dealing here with a double expansion.  An expansion in the
coefficient $\kappa$ (appearing in the Dirac field action) was needed to
define quark loops on the lattice; the sum of the $\kappa$ expansion is a
sum over all possible quark loops.  The $g^{-2}$ expansion is needed to
define surfaces filling in the quark loops.  The sum of the $g$ expansion
is a sum over all such surfaces.  This is precisely the structure
appearing in string models:  combined sums over quark loops and
interpolating surfaces.  However, the loops and surfaces of the
gauge-field theory are defined on a lattice whereas the loops and
surfaces of the string models are defined on a continuum.  It may not be
easy to derive quantitative relations between the two types of surfaces.
